% =====================================================================
% 第1章 Git編 — 自分のPCの中でGitを使う
% 元原稿: 01_Git編.md
% =====================================================================
\chapter{Git編 —自分のPCの中でGitを使う—}

この章では、\textbf{自分のPCの中だけで完結するGitの基本}を説明します。GitHub（クラウド）とつなぐ話は第2章で扱います。まずはここで「変更を記録する」感覚をつかみましょう。

\begin{hintbox}
\textbf{この章の読み方}
\begin{itemize}[leftmargin=1.4em, nosep]
  \item とりあえず動かしたい人 → 「インストール」→「最初の設定」→「【実践】はじめてのコミット」だけ写経すればOK
  \item コマンドをド忘れした人 → 『逆引きチートシート』へ
  \item エラーが出て困った人 → 『トラブルシューティング』へ
\end{itemize}
\end{hintbox}

% ---------------------------------------------------------------------
\section{そもそもGitとは}

Gitは\textbf{ファイルの変更履歴を記録・管理するツール}です。一言でいうと「\textbf{いつでも戻れるセーブポイントを作れる仕組み}」です。

たとえば、レポートをこんなふうに管理していませんか。

\begin{codebox}[text]
\begin{code}
レポート.docx
レポート_最終.docx
レポート_最終2.docx
レポート_本当に最終.docx
\end{code}
\end{codebox}

このやり方だと、次のことに困ります。

\begin{itemize}[leftmargin=1.4em]
  \item \textbf{どれが本当に最新か分からない}
  \item \textbf{ファイルが増え続けて散らかる}
  \item \textbf{いつ・誰が・どこを・なぜ変えたのか分からない}
\end{itemize}

コピーを取っておけば「昔の状態に戻す」ことはできます。でも「\textbf{どこが違うのか}」「\textbf{なぜ変えたのか}」までは分かりません。ここがGitとの決定的な違いです。

Gitを使うと、\textbf{ファイルは1つのまま}、変更の節目ごとに「\textbf{コミット}」という記録（セーブポイント）を残せます。各コミットには「いつ・誰が・何を・なぜ変えたか」に加えて、\textbf{前回との差分}まで残るので、後からいつでも正確にたどれます。

\begin{numfig}{コミットを積み重ねると、変更履歴が1本の線としてつながる}{fig:commit-chain}
\begin{tikzpicture}[node distance=1.1cm]
  \node[flownode, fill=gitBlue] (A) {昔のレポート\\（コミットA）};
  \node[flownode, fill=gitBlue, right=of A] (B) {少し進めた\\（コミットB）};
  \node[flownode, fill=gitGreen, right=of B] (C) {完成\\（コミットC）};
  \draw[flowarrow] (A) -- (B);
  \draw[flowarrow] (B) -- (C);
\end{tikzpicture}
\end{numfig}

% ---------------------------------------------------------------------
\section{Gitの3つのエリア（最重要の概念）}

Gitでは、ファイルは次の3つのエリアを通って記録されます。\textbf{ここだけは必ず押さえてください。}

\begin{hintbox}
下の図に出てくる \cmd{git add} と \cmd{git commit} は、この後の「実践」で実際に使うコマンドです。ここでは\textbf{エリアの間を移す操作の名前}とだけ捉えてください。
\end{hintbox}

3つのエリアは、\textbf{机 → 段ボール箱 → 倉庫}にたとえると分かりやすいです（図\ref{fig:areas-concept}）。作業ツリーは原本を広げた\textbf{作業机}、ステージングは変更を詰める\textbf{段ボール箱}、ローカルリポジトリは箱を保管する\textbf{倉庫}です。

\figimage[0.95]{01_git/git_area.png}{3つのエリアの概念図（机 → 段ボール箱 → 倉庫）}{fig:areas-concept}

この流れをコマンドで表すと、\cmd{git add} と \cmd{git commit} で次のエリアへ送っていくことになります（図\ref{fig:areas-flow}）。

\begin{numfig}{3つのエリアと、それらをつなぐ \texttt{git add} / \texttt{git commit}}{fig:areas-flow}
\begin{tikzpicture}
  \node[flownode, fill=gitRed, text width=2.4cm] (W)
        {作業ツリー\\{\scriptsize (Working Tree)}\\原本を広げた作業机};
  \node[flownode, fill=gitOrange, text width=2.4cm, right=3.6cm of W] (S)
        {ステージング\\{\scriptsize (Staging Area)}\\写しを詰める段ボール箱};
  \node[flownode, fill=gitGreen, text width=2.4cm, right=3.6cm of S] (R)
        {ローカルリポジトリ\\{\scriptsize (Local Repository)}\\写しを保管する倉庫};
  \draw[flowarrow] (W) -- (S) node[edgelabel, midway, above] {git add};
  \draw[flowarrow] (S) -- (R) node[edgelabel, midway, above] {git commit};
\end{tikzpicture}
\end{numfig}

\begin{center}
\begin{tabularx}{\linewidth}{@{}l X X@{}}
\toprule
エリア & これは何か & たとえると \\
\midrule
\textbf{作業ツリー} & 実際にファイルを編集している、いつものフォルダ & 原本を広げた作業机 \\
\textbf{ステージング} & 「次のコミットに含める変更」を詰めるエリア & コミット1回分の段ボール箱 \\
\textbf{ローカルリポジトリ} & コミットとして履歴が保存されるエリア（\cmd{.git} の中） & 箱を保管する倉庫 \\
\bottomrule
\end{tabularx}
\end{center}

大事なのは、\cmd{git add} も \cmd{git commit} も「\textbf{写しを取って次のエリアに置く}」操作だという点です。元のエリアには元のまま残るので、先に進めても空になりません。

\begin{itemize}[leftmargin=1.4em]
  \item \cmd{git add}：机の上の原本の写しを箱に入れる（原本は机に残る）
  \item \cmd{git commit}：箱の中身の写しを倉庫に保管する（箱の中身も残る）
\end{itemize}

\begin{notebox}
Gitは各エリアが「その時点の写し（スナップショット）」を持つように動きます。だからコミットは\textbf{差分ではなくスナップショット}として記録されます。
\end{notebox}

つまり基本操作は \textbf{「編集する → \cmd{add} → \cmd{commit}」} の繰り返しです。本書ではこれを\textbf{黄金フロー}と呼びます。この3ステップが分かれば、Gitの8割は理解できたようなものです。

% ---------------------------------------------------------------------
\section{Gitのインストール}

まず、Gitが\textbf{すでに入っていないか}を確認しましょう。ターミナル（WindowsはPowerShellやコマンドプロンプト）で次を実行します。

\begin{codebox}[bash]
\begin{code}
git --version
# 出力例: git version 2.45.0
\end{code}
\end{codebox}

\begin{hintbox}
上のように\textbf{バージョンが表示されたら、Gitはすでにインストール済み}です。以下のOS別のインストール手順は\textbf{飛ばして}、次の「最初の設定」へ進んでください。\cmd{command not found}（見つからない）と出た場合だけ、お使いのOSの手順に従ってインストールします。インストール後にもう一度 \cmd{git --version} を実行し、バージョンが表示されれば成功です。
\end{hintbox}

\begin{warnbox}
\textbf{インストール直後は、開いていたターミナルではまだ \cmd{git --version} が反映されないことがあります。} PATH（コマンドの置き場所を示す設定）がターミナルの起動時に読み込まれるためです。\cmd{command not found} のままなら、\textbf{ターミナルをいったん閉じて開き直して}から、もう一度実行してください。
\end{warnbox}

\subsection{Windows}

Windowsでは「Git for Windows」をインストールします。これには \textbf{Git Bash} という、macOS/Linux風のコマンドが使えるターミナルが付いてきます。本書のWindowsの説明はこの Git Bash を使う前提です。

\begin{enumerate}[leftmargin=1.6em]
  \item \url{https://git-scm.com/download/win} からインストーラをダウンロードし、ダブルクリックで起動します。
  \item インストーラの画面が続きます。\textbf{基本はすべてデフォルトのまま「Next」で進めて問題ありません。} ただし次の3つの画面だけは選択を確認してください。
\end{enumerate}

\noindent\textbf{① 既定のエディタ}

デフォルトはVimですが不慣れだと難しいため、VS Codeを使う予定なら \textbf{「Use Visual Studio Code as Git's default editor」} を選びます。
\screenshot[0.62]{01_git/win_install_05.png}{既定エディタの選択}

\noindent\textbf{② 初期ブランチ名}

\textbf{「Override the default branch name…」} を選び、\cmd{main} と入力します（世界的に \cmd{main} が標準）。
\screenshot[0.62]{01_git/win_install_06.png}{初期ブランチ名の設定}

\noindent\textbf{③ 改行コードの扱い}

WindowsとほかのOSで改行コードが異なるため重要です。複数環境で開発するなら \textbf{「Checkout as-is, commit Unix-style line endings」} を選びます。
\screenshot[0.62]{01_git/win_install_10.png}{改行コードの設定}

残りの画面はデフォルトのまま進め、最後に「Install」をクリックします。インストール後、スタートメニューから \textbf{Git Bash} を起動し、\cmd{git --version} でバージョンが出れば完了です。

\subsection{macOS}

次のどちらか一方でOKです。

\noindent\textbf{方法A: Xcode Command Line Tools（最も手軽）}

ターミナル（\cmd{アプリケーション > ユーティリティ > ターミナル.app}）を開き、次を実行します。

\begin{codebox}[bash]
\begin{code}
xcode-select --install
\end{code}
\end{codebox}

ダイアログが出たら「インストール」をクリックします。

\noindent\textbf{方法B: Homebrew（最新版を使いたい場合）}

Homebrew（\url{https://brew.sh/index_ja}）が入っていれば、次で最新のGitが入ります。

\begin{codebox}[bash]
\begin{code}
brew install git
\end{code}
\end{codebox}

どちらの場合も、\cmd{git --version} でバージョンが表示されれば完了です。

\subsection{Linux}

お使いのディストリビューションのパッケージマネージャを使います。自分の環境に合うものを1つ選んで実行します。

\begin{codebox}[bash]
\begin{code}
# Debian / Ubuntu 系
sudo apt update && sudo apt install git

# Fedora / RHEL 系
sudo dnf install git

# Arch 系
sudo pacman -S git
\end{code}
\end{codebox}

\cmd{git --version} でバージョンが表示されれば完了です。

% ---------------------------------------------------------------------
\section{最初の設定}

インストール直後に\textbf{一度だけ}行う設定です。コミットに記録される「あなたの名前とメールアドレス」を登録します。ここで設定した名前・メールが、後でGitHub上の履歴にも表示されます。

ターミナル（Windowsは Git Bash）を開き、次を順に実行します。

\begin{codebox}[bash]
\begin{code}
# あなたの名前を登録（GitHubのアカウント名に合わせるのがおすすめ）
git config --global user.name "your-name"

# あなたのメールアドレスを登録（GitHubで使う予定のメール）
git config --global user.email "you@example.com"

# 設定内容を確認する
git config --global --list
\end{code}
\end{codebox}

\begin{warnbox}
名前やメールにスペースが含まれる場合は、上のように \textbf{ダブルクォート \cmd{"} で囲んでください}。
\end{warnbox}

\cmd{--global} は「このPC全体の共通設定」という意味です。一度やれば、以降すべてのリポジトリに適用されます。

\begin{hintbox}
おすすめの追加設定（新規リポジトリの初期ブランチ名を \cmd{main} に統一）
\begin{codebox}[bash]
\begin{code}
git config --global init.defaultBranch main
\end{code}
\end{codebox}
\end{hintbox}

% ---------------------------------------------------------------------
\section{【実践】はじめてのコミット}

ここが本番です。\textbf{黄金フロー（編集 → \cmd{add} → \cmd{commit}）}を、実際に手を動かして体験します。\textbf{この通りに写経すれば、最初のコミットが作れます。}

\subsection{ステップ0：作業フォルダを用意して移動する}

\begin{codebox}[bash]
\begin{code}
# 練習用フォルダを作って中に入る
mkdir git-practice
cd git-practice
\end{code}
\end{codebox}

\begin{hintbox}
\cmd{cd}（change directory）はフォルダを移動するコマンドです。以降の操作は「今いるフォルダ」に対して行われます。
\end{hintbox}

\subsection{ステップ1：リポジトリを作る（git init）}

\begin{codebox}[bash]
\begin{code}
git init
\end{code}
\end{codebox}

これで今のフォルダに \cmd{.git} という隠しフォルダが作られ、Gitでの管理が始まります。\textbf{これが、あなたのPCの中の「ローカルリポジトリ」}です。履歴（コミット）はこの \cmd{.git} の中に保存されます（中身を直接触る必要はありません）。

\begin{codebox}[bash]
\begin{code}
# git initに成功すると、次のように表示されます
hoge@n171h100 git-practice % git init
Initialized empty Git repository in /Users/hoge/git-practice/.git/
\end{code}
\end{codebox}

\screenshot[0.62]{01_git/01_git_init_folder.png}{git init後の.gitフォルダ（隠しファイル表示をONにすると見えます）}

\subsection{ステップ2：ファイルを作る}

試しにテキストファイルを1つ作ります。

\begin{codebox}[bash]
\begin{code}
echo "Hello Git" > hello.txt
\end{code}
\end{codebox}

\subsection{ステップ3：変更を確認する（git status）}

\begin{codebox}[bash]
\begin{code}
git status
\end{code}
\end{codebox}

\cmd{hello.txt} が「Untracked files（まだ管理されていないファイル）」として赤色で表示されます。
\screenshot[0.62]{01_git/01_git_status.png}{git statusの出力}

\subsection{ステップ4：ステージングに追加する（git add）}

\begin{codebox}[bash]
\begin{code}
git add hello.txt
\end{code}
\end{codebox}

これで \cmd{hello.txt} が「次のコミットに含める変更」として選ばれました。もう一度 \cmd{git status} を実行すると、今度は緑色で表示されます。
\screenshot[0.62]{01_git/01_git_add_status.png}{git add後のgit statusの出力}

\subsection{ステップ5：コミットする（git commit）}

\begin{codebox}[bash]
\begin{code}
git commit -m "最初のコミット"
\end{code}
\end{codebox}

\cmd{-m} の後ろの文字列が\textbf{コミットメッセージ}（変更内容の一言メモ）です。これでコミット（セーブポイント）が1つ作られました。

\subsection{ステップ6：履歴を確認する（git log）}

\begin{codebox}[bash]
\begin{code}
git log
\end{code}
\end{codebox}

\screenshot[0.62]{01_git/01_git_commit_log.png}{git logの出力}

今作ったコミットが、作成者・日時・メッセージとともに表示されます。\cmd{q} キーで表示を終了します。

\subsection{黄金フローのまとめ}

上の\textbf{ステップ4（\cmd{git add}）→ ステップ5（\cmd{git commit}）}が黄金フローの中核です。図にすると、図\ref{fig:golden-flow}のように3つの繰り返しになります。

\begin{numfig}{黄金フロー：編集 → \texttt{git add} → \texttt{git commit} の繰り返し}{fig:golden-flow}
\begin{tikzpicture}[scale=1.5, transform shape]
  \node[flownode, fill=gitRed] (E) {① ファイルを編集};
  \node[flownode, fill=gitOrange, right=3.4cm of E] (S) {ステージング};
  \node[flownode, fill=gitGreen, right=3.4cm of S] (R) {コミットとして\\記録};
  \draw[flowarrow] (E) -- (S) node[edgelabel, midway, above] {② git add};
  \draw[flowarrow] (S) -- (R) node[edgelabel, midway, above] {③ git commit};
  \draw[flowarrow] (R) to[out=-30, in=-150]
        node[edgelabel, midway, below] {繰り返す} (E);
\end{tikzpicture}
\end{numfig}

実際の作業は、この \textbf{「編集 → add → commit」} をひたすら繰り返すだけです。

\begin{codebox}[bash]
\begin{code}
# 日常的に使う最小セット
git add .                       # 変更した全ファイルをステージングに追加
git commit -m "変更内容のメモ"   # コミットとして記録
\end{code}
\end{codebox}

\begin{hintbox}
\cmd{git add .} の \cmd{.}（ドット）は「今いるフォルダ以下のすべての変更」という意味です。
\end{hintbox}

\subsection{いつ \texttt{add} / \texttt{commit} すればいい？}

最初に迷うのが\textbf{コミットするタイミング}です。目安はシンプルで、\textbf{「1つの作業に区切りがついたら」コミットする}——これだけです。たとえば次のような節目です。

\begin{itemize}[leftmargin=1.4em]
  \item 機能を1つ作り終えた
  \item 文章を1節書けた／図を1つ追加した
  \item バグを1つ直した
  \item 設定を変えて、動くことを確認できた
\end{itemize}

ポイントは、\textbf{後から見て「何をしたか」を一言で言えるまとまり}にすることです。その一言が、そのまま\textbf{コミットメッセージ}（例：\cmd{ログイン画面を追加}）になります。

\begin{hintbox}
\textbf{こまめに、小さくコミットするのがおすすめ}です。小さく区切っておくと、\textbf{後で戻したいときにその1コミットだけ取り消せ}、\textbf{不具合が入った箇所も見つけやすく}なります。逆に何日分もためて1回でコミットすると、後から分けられず、メッセージも「いろいろ修正」のように曖昧になりがちです。
\end{hintbox}

\textbf{\cmd{add} と \cmd{commit} の役割の違い}も、このタイミングで押さえておきましょう。

\begin{itemize}[leftmargin=1.4em]
  \item \cmd{git add}：これから作るコミットに\textbf{含める変更を選ぶ}（コミット直前の準備）
  \item \cmd{git commit}：選んだ変更で\textbf{セーブポイントを確定する}（区切りをつける）
\end{itemize}

ふだんは \cmd{git add .} → \cmd{git commit} を\textbf{続けて実行すればOK}です。「変更した中の一部だけをコミットしたい」ときにだけ、\cmd{add} で\textbf{選ぶ}意味が出てきます。

% ---------------------------------------------------------------------
\section{状態を確認するコマンド}

Gitは「今どういう状態か」をこまめに確認しながら使うのが安全です。困ったらまず \cmd{git status} を打つ癖をつけましょう。

\begin{center}
\begin{tabularx}{\linewidth}{@{}l X@{}}
\toprule
コマンド & 何が分かるか \\
\midrule
\cmd{git status} & どのファイルが変更され、どのエリアにあるか \\
\cmd{git log} & これまでのコミット履歴（作成者・日時・メッセージ） \\
\cmd{git log -{}-oneline} & コミット履歴を1行ずつ簡潔に表示 \\
\cmd{git diff} & まだステージングしていない変更の中身（どこを書き換えたか） \\
\cmd{git diff -{}-staged} & ステージング済みの変更の中身 \\
\bottomrule
\end{tabularx}
\end{center}

このうち \textbf{\cmd{git diff}} は「どこをどう書き換えたか」を確認するのによく使います。追加された行が \cmd{+} 付きの緑、削除された行が \cmd{-} 付きの赤で表示されます。

\begin{notebox}
\cmd{git add} でステージングした後は \cmd{git diff} に何も出なくなります。そのときは \cmd{git diff -{}-staged} を使います。さらに \cmd{git commit} した後は、\cmd{git log} で確認します。
\end{notebox}
\screenshot[0.92]{01_git/01_git_log_diff.png}{git logとgit diffの出力}

% ---------------------------------------------------------------------
\section{.gitignore：記録したくないファイルを除外する}

プロジェクトには「Gitで記録したくないファイル」があります。たとえば次のようなものです。

\begin{itemize}[leftmargin=1.4em]
  \item パスワードやAPIキーなどの\textbf{秘密情報}（誤って公開すると危険）
  \item 自動生成される一時ファイル（\cmd{.DS\_Store}、\cmd{\_\_pycache\_\_/}、LaTeXの \cmd{.aux} など。作り直せるので記録する意味がない）
  \item サイズの大きいデータやビルド成果物（リポジトリが重くなる。GitHubには容量の上限もある）
\end{itemize}

これらは \cmd{.gitignore} というファイルに書いておくと、Gitが\textbf{最初から無視}してくれます（\cmd{git status} にも出てきません）。リポジトリの一番上の階層に \cmd{.gitignore} を作り、無視したいものを1行ずつ書きます。

\begin{codebox}[text]
\begin{code}
# macOS が作る隠しファイル
.DS_Store

# Python のキャッシュ
__pycache__/
*.pyc

# LaTeX の中間ファイル
*.aux
*.log
*.out

# 秘密情報（絶対にコミットしない）
.env
secrets.txt
\end{code}
\end{codebox}

\begin{warnbox}
\textbf{一度コミットしてしまったファイルは、後から \cmd{.gitignore} に書いても無視されません。} 秘密情報は\textbf{最初のコミット前}に \cmd{.gitignore} へ登録してください。
\end{warnbox}

\begin{hintbox}
言語やフレームワークごとの \cmd{.gitignore} テンプレートは \url{https://github.com/github/gitignore} で、自動生成は \url{https://www.toptal.com/developers/gitignore} で手に入ります。
\end{hintbox}

% ---------------------------------------------------------------------
\section*{この章のまとめ}
\addcontentsline{toc}{section}{この章のまとめ}

\begin{itemize}[leftmargin=1.4em]
  \item Gitは「いつでも戻れるセーブポイント」を作る変更履歴管理ツール
  \item ファイルは \textbf{作業ツリー → ステージング → ローカルリポジトリ} の3エリアへ「写し」が渡っていく（元のエリアは空にならない）
  \item 基本操作は \textbf{「編集 → \cmd{git add} → \cmd{git commit}」} の繰り返し
  \item 困ったら \cmd{git status} で現在地を確認する
\end{itemize}

次は、ここで作った履歴を\textbf{クラウド（GitHub）に上げて共有する}方法です（第2章）。
